\chapter{الجمع بين الأحاديث، ودلالة النهي، وأقسام العلم}

\section{الجمع بين الأحاديث (المزابنة والعرايا)}
الحمد لله والصلاة والسلام على رسول الله، أما بعد؛ فيواصل الإمام الشافعي رحمه الله تعالى بيان الأحاديث التي ظاهرها الاختلاف وكيفية الجمع بينها.
ومن ذلك: ما روي عن ابن عمر رضي الله عنهما أن رسول الله ﷺ «نهى عن المزابنة». والمزابنة: بيع الثمر بالتمر كيلاً (أي بيع الرطب على النخل بالتمر الجاف). وفي حديث سعد بن أبي وقاص رضي الله عنه أنه سُئل النبي ﷺ عن بيع التمر بالرطب، فقال: «أينقص الرطب إذا يبس؟»، قالوا: نعم، فنهى عن ذلك. 
فهذه الأحاديث تدل على تحريم هذا البيع صراحة لوجود التفاضل والجهالة فيه.

وفي المقابل، روي عن زيد بن ثابت رضي الله عنه أن رسول الله ﷺ «رخص لصاحب العرية أن يبيعها بخرصها». والعرايا: أن يمتلك الفقير رطباً على النخل ويحتاج إلى تمر ليأكله، فرخص له النبي ﷺ أن يبيع هذا الرطب بتمر يخرصه (يقدره) أهل الخبرة، ما لم يبلغ خمسة أوسق.
ظاهر هذين الحديثين التعارض؛ فالأول ينهى مطلقاً، والثاني يبيح. والجمع بينهما عند الشافعي: أن النهي عام عن المزابنة لحرمة الجهالة والغرر والربا، وأن العرايا رخصة مستثناة من هذا النهي العام لحاجة الفقراء بحدود وضوابط (وهي الخرص وعدم تجاوز خمسة أوسق). فالعام باقٍ على عمومه، والعرايا مخصوصة ومستثناة للحاجة.

\section{الجمع بين بيع المجهول وإباحة السَّلَم}
ومن ذلك أيضاً: حديث حكيم بن حزام رضي الله عنه، حيث قال: «يا رسول الله يأتيني الرجل فيريد مني المبيع ليس عندي، أفأبتاعه له من السوق؟» فقال ﷺ: «لا تبع ما ليس عندك».
ويقابله حديث ابن عباس رضي الله عنهما: «قدم رسول الله ﷺ المدينة وهم يسلفون في التمر السنتين والثلاث، فقال: من أسلف في شيء فليُسلف في كيل معلوم، ووزن معلوم، إلى أجل معلوم». والسَّلَم (السَّلَف) هو بيع شيء ليس موجوداً عند العقد.

\textbf{وجه الجمع بينهما:}
النهي في حديث حكيم عن بيع «عين» ليس يملكها البائع وليست بحضرته وقد تهلك قبل تسليمها (كأن يبيع سيارة محددة لا يملكها). أما «السَّلَم» المباح، فهو بيع «موصوف في الذمة» مضمون على البائع (كأن يبيع طناً من القمح بصفة كذا يسلمه بعد سنة)، فهذا مباح تيسيراً على الناس. فلا تعارض بين النهي عن بيع المعيَّن الغائب، وبين إباحة بيع الموصوف المضمون في الذمة. 
فالجمع بين الأحاديث هو مسلك العلماء، ولا يصار إلى دعوى التعارض والنسخ إلا عند تعذر الجمع.

\section{صفة نهي الله ورسوله (بين التحريم والكراهة)}
قسم الشافعي رحمه الله ما نهى عنه الشرع إلى قسمين: نهي يقتضي التحريم والفساد، ونهي يقتضي الكراهة والتنزيه (أو الإرشاد).

\subsection{النهي الذي يقتضي التحريم والفساد}
إذا جاء النهي عن شيء في أصل العقود أو العبادات، فإنه يقتضي التحريم وإبطال العقد. كنهي الله ورسوله عن نكاح المرأة على عمتها أو خالتها، ونهيه عن الجمع بين الأختين، ونهيه عن نكاح المتعة، ونهيه عن نكاح الشغار (وهو أن يزوجه أخته على أن يزوجه الآخر أخته دون مهر).
فأي نكاح أو بيع وقع على هذه الوجوه المنهي عنها فهو مفسوخ وباطل.
ويقرر الشافعي قاعدة عظيمة هنا: أن «المعصية لا تُحِلُّ مُحَرَّماً». فالأصل في الأبضاع والأموال التحريم، ولا تستباح إلا بما شرعه الله. والغاية المشروعة لا تُنال بوسيلة محرمة؛ فالتمكين للدين لا ينال بالانحراف عن منهج السلف واللجوء للبدع أو النظم المستوردة، وما عند الله لا ينال بمعصيته.

\subsection{النهي الذي يقتضي الكراهة والتنزيه (الإرشاد)}
قد يأتي النهي عن شيء أصله مباح، ويكون النهي لعلة أخرى خارجة عن ذات الشيء، فهذا نهي أدب وإرشاد. ومن أمثلته:
\begin{itemize}
    \item \textbf{النهي عن اشتمال الصمَّاء والاحتباء:} الأصل في الثياب الإباحة، لكن نهى النبي ﷺ أن يلتف الرجل بثوب واحد لا يستطيع إخراج يده منه (الصمَّاء)، ونهى أن يحتبي (يضم ركبتيه إلى صدره) في ثوب واحد ليس تحته غيره. والنهي هنا ليس لكون الثوب محرماً، بل خشية انكشاف العورة، فهو نهي إرشاد.
    \item \textbf{النهي عن الأكل من وسط الصحفة:} الأصل إباحة الطعام، لكن نهى النبي ﷺ عن الأكل من أعلى الصحفة (وسطها) وأمر بالأكل من الحواف؛ لأن البركة تنزل في وسطها. فهذا نهي تأديب وليس تحريماً لذات الطعام.
    \item \textbf{النهي عن التعريس بالطريق:} التعريس هو النزول للراحة والنوم في السفر. وقد نهى النبي ﷺ عن التعريس في قارعة الطريق؛ لأنها ممر الهوام (كالحيات) ومسلك الدواب. فالنهي هنا للإرشاد والمحافظة على المسلم، وليس تحريماً لجلوسه على الأرض.
\end{itemize}
ومن فعل شيئاً من هذا عالماً بالنهي فهو عاصٍ (لمخالفته الأمر)، ولكنه ليس كمن فعل المحرمات القطعية المفسدة للعقود، لأن هذه تتعلق بأمور مباحة الأصل.

\section{أقسام العلم الشرعي}
عقد الشافعي باباً نفيساً في أقسام العلم، فبين أن العلم ينقسم إلى قسمين:

\subsection{علم العامة (فرض العين)}
وهو العلم الذي لا يسع مكلفاً بالغاً عاقلاً جهله. وهو نص كتاب الله وما تواتر من سنة رسول الله ﷺ في أصول الدين، مثل: فرضية الصلوات الخمس، وصيام رمضان، وحج البيت، وزكاة الأموال، وتحريم الزنا والقتل والسرقة والخمر. 
هذا العلم يعلمه العامة عن العامة، ولا يُقبل فيه التأويل، ولا يُعذر فيه أحد بالجهل، ولا يختلف فيه مسلمان.

\subsection{علم الخاصة (فرض الكفاية)}
وهو العلم بتفاصيل الأحكام، وفروع المسائل، ودقائق السنن والناسخ والمنسوخ، وما يستنبط بالقياس والاجتهاد.
هذا النوع من العلم ليس فرضاً على جميع الأمة، بل هو (فرض كفاية)؛ إذا قام به مَن فيهم غَناء وكفاية من طلبة العلم والعلماء، سقط الإثم عن الباقين، ونال القائمون به الفضل والدرجات العُلى.

واستدل الشافعي على فرض الكفاية بآيات الجهاد؛ فإن الله أوجب الجهاد، ولكنه رفع الحرج عن القاعدين (من غير أولي الضرر) ووعدهم الحسنى، وفضل المجاهدين عليهم درجات. فلو كان فرض عين لأثم القاعدون. وكذلك قال تعالى: ﴿وَمَا كَانَ الْمُؤْمِنُونَ لِيَنفِرُوا كَافَّةً ۚ فَلَوْلَا نَفَرَ مِن كُلِّ فِرْقَةٍ مِّنْهُمْ طَائِفَةٌ لِّيَتَفَقَّهُوا فِي الدِّينِ﴾، فجعل النفير لطلب العلم كفاية في طائفة تنذر قومها إذا رجعوا إليهم.
ومن أمثلة فرض الكفاية أيضاً: صلاة الجنازة وتغسيل الميت ودفنه، ورد السلام؛ إذا قام به واحد من القوم أجزأ عن الباقين.